\documentclass[11pt]{article}

\usepackage[T1]{fontenc}
\usepackage[utf8]{inputenc}
\usepackage{lmodern}
\usepackage{microtype}
\usepackage[a4paper,margin=30mm,headheight=14pt]{geometry}
\usepackage{amsmath,amssymb,amsthm,mathtools,bm,mathrsfs}
\usepackage{aliascnt}
\usepackage{booktabs}
\usepackage{enumitem}
\usepackage{fancyhdr}
\usepackage{needspace}
\usepackage{xcolor}
\usepackage{hyperref}
\usepackage[nameinlink,noabbrev]{cleveref}

\definecolor{linkblue}{RGB}{32,72,118}
\hypersetup{
  colorlinks=true,
  linkcolor=linkblue,
  citecolor=linkblue,
  urlcolor=linkblue,
  pdftitle={One Fluid History Through the Terminal Time},
  pdfauthor={Thomas Birnie}
}

\allowdisplaybreaks[2]
\setlength{\parindent}{1.35em}
\setlength{\parskip}{0.18em}
\setlist{itemsep=0.25em,topsep=0.45em}

\pagestyle{fancy}
\fancyhf{}
\fancyhead[L]{\small One fluid history through the terminal time}
\fancyhead[R]{\small Thomas Birnie}
\fancyfoot[C]{\thepage}

\newtheorem{theorem}{Theorem}[section]
\newaliascnt{proposition}{theorem}
\newtheorem{proposition}[proposition]{Proposition}
\aliascntresetthe{proposition}
\newaliascnt{lemma}{theorem}
\newtheorem{lemma}[lemma]{Lemma}
\aliascntresetthe{lemma}
\newaliascnt{corollary}{theorem}
\newtheorem{corollary}[corollary]{Corollary}
\aliascntresetthe{corollary}
\theoremstyle{definition}
\newaliascnt{definition}{theorem}
\newtheorem{definition}[definition]{Definition}
\aliascntresetthe{definition}
\theoremstyle{remark}
\newaliascnt{remark}{theorem}
\newtheorem{remark}[remark]{Remark}
\aliascntresetthe{remark}
\crefname{theorem}{theorem}{theorems}
\crefname{proposition}{proposition}{propositions}
\crefname{lemma}{lemma}{lemmas}
\crefname{corollary}{corollary}{corollaries}
\crefname{definition}{definition}{definitions}
\crefname{remark}{remark}{remarks}

\newcommand{\R}{\mathbb R}
\newcommand{\N}{\mathbb N}
\newcommand{\Nzero}{\mathbb N_0}
\newcommand{\Hs}{H^\infty}
\newcommand{\HsSigma}{H^\infty_\sigma}
\newcommand{\Schwartz}{\mathcal S}
\newcommand{\Leray}{\mathbb P}
\newcommand{\cR}{\mathcal R}
\newcommand{\Szero}{\mathscr S_0}
\newcommand{\norm}[1]{\left\lVert #1\right\rVert}
\newcommand{\pair}[2]{\left\langle #1,#2\right\rangle}
\newcommand{\dd}{\,\mathrm d}
\newcommand{\e}{\mathrm e}
\newcommand{\diver}{\operatorname{div}}
\DeclareMathOperator*{\esssup}{ess\,sup}

\title{\bfseries One Fluid History Through the Terminal Time\\[0.35em]
\large Global Closure from Compatible Mixed Sobolev Rectangles for the
Three-Dimensional Navier--Stokes Equations}
\author{Thomas Birnie}
\date{}

\begin{document}
\maketitle

\begin{abstract}
For smooth divergence-free finite-energy data on \(\R^3\), the
datum-generated whole-terminal theorem supplies every finite adjacent
Sobolev rectangle of the one maximal Navier--Stokes history.  We prove the
global closure of that theorem.  Direct differentiation produces the
pressure-normalized mixed jet.  Different finite parts agree because they are
derivatives of the same solution, and their overlap gives its complete
spatial row, not an infinite norm sum or a sequence of replacement solutions.
The equation then makes the nonlinear pressure-completed interaction
integrable at every Sobolev order.  The velocity reaches the alleged terminal
time in every \(H^m\), and the compatible traces define one endpoint in
\(H^\infty\).  Adjacent rows carry every higher temporal response to that
same time.  The differentiated pressure and momentum equations identify the
actual endpoint jet.  A uniform late-slice local construction overlaps the
original history, and uniqueness extends it beyond the alleged maximal time.
This proves global smoothness from the datum theorem.  The critical
\(H^{3/2}\) dissipation and \(H^{1/2}\) height follow as consequences of the
same compatible family.
\end{abstract}

\section{The terminal problem belongs to one history}

Fix \(\nu>0\) and a real-valued divergence-free datum
\[
 u_0\in \Schwartz_\sigma(\R^3).
\]
We consider the unforced incompressible Navier--Stokes equations
\begin{equation}
 \partial_tu+(u\cdot\nabla)u+\nabla p=\nu\Delta u,
 \qquad \nabla\cdot u=0,
 \qquad u(0)=u_0.
 \label{eq:NS}
\end{equation}
The pressure is normalized by decay at spatial infinity.  Classical local
theory gives a unique smooth solution on a maximal interval
\([0,T_*)\), where \(0<T_*\leq\infty\); see
\cite{FujitaKato,Kato}.  Global regularity is the assertion that this maximal
time cannot be finite in the whole-space Millennium formulation
\cite{Fefferman}.

The point of departure is simple but decisive.  If \(T_*<\infty\), the object
approaching \(T_*\) is still the velocity field generated by \(u_0\).  Its
time derivatives, spatial derivatives, pressure responses, and endpoint
traces are not rival evolutions.  They are simultaneous readings of one
motion.  The proof succeeds only if every finite reading reaches the same
terminal state.

The ordinary energy calculation already preserves part of that identity.
Taking the \(L^2\) inner product of \eqref{eq:NS} with \(u\), using
incompressibility, and integrating in time gives
\begin{equation}
 \frac12\norm{u(t)}_2^2
 +\nu\int_0^t\norm{\nabla u(s)}_2^2\dd s
 =\frac12\norm{u_0}_2^2
 \qquad (0<t<T_*).
 \label{eq:energy}
\end{equation}
Thus kinetic energy cannot be the quantity that disappears at a finite
endpoint.  What remains is a derivative-level problem: can all finite
Sobolev readings of the same solution be carried through the time at which
classicality was alleged to stop?

The whole-terminal datum theorem answers this question through finite adjacent
rectangles.  Before giving the definitions, we state the closure it forces.

\begin{theorem}[Global closure of the datum theorem]
\label{thm:global}
Let \(u_0\in\Schwartz_\sigma(\R^3)\) and \(\nu>0\).  The unique maximal
classical solution of \eqref{eq:NS} obeying the whole-terminal datum theorem
in \cref{thm:rectangles} satisfies \(T_*=\infty\),
\[
 u,p\in C^\infty([0,\infty)\times\R^3),
\]
and the energy identity \eqref{eq:energy} holds for every finite time.
\end{theorem}

Assume \(T_*<\infty\).  The energy identity supplies only the base control
\[
 u\in L^\infty(0,T_*;L^2)\cap L^2(0,T_*;H^1).
\]
It does not place a chosen higher response and its time derivative in an
adjacent Sobolev pair uniformly through \(T_*\).  Such a pair is exactly what
creates a middle-order trace.  The datum theorem in
\cref{sec:rectangles} supplies every fixed finite pair on the whole terminal
interval.  Their shared entries are derivatives of the same history, so the
complete spatial row reaches one \(H^\infty\) endpoint.  The equation
recovers the actual pressure-complete jet there, and a late-slice restart
continues the history past \(T_*\).  No bound uniform over all derivative
orders is used.

\section{The equation generates its pressure-complete jet}
\label{sec:jet}

Taking the divergence of \eqref{eq:NS} and using \(\nabla\cdot u=0\) gives
\begin{equation}
 -\Delta p=\partial_i\partial_j(u_i u_j),
 \qquad p=R_iR_j(u_i u_j).
 \label{eq:pressure}
\end{equation}
Here \(R_i\) denotes the \(i\)-th Riesz transform.  The second identity fixes
the chosen decaying pressure and makes its dependence on the whole velocity
field explicit.  Pressure is a nonlocal simultaneous response of the same
history, not an additional datum.

Set
\[
 V_n:=\partial_t^n u,
 \qquad Q_n:=\partial_t^n p.
\]
Leibniz' rule and \eqref{eq:pressure} give, for every \(n\in\Nzero\),
\begin{align}
 Q_n
 &=R_iR_j\sum_{a=0}^{n}\binom na V_{a,i}V_{n-a,j},
 \label{eq:pressure-recurrence}\\
 V_{n+1}
 &=\nu\Delta V_n
 -\sum_{a=0}^{n}\binom na(V_a\cdot\nabla)V_{n-a}
 -\nabla Q_n.
 \label{eq:velocity-recurrence}
\end{align}
The momentum recurrence advances the velocity response by one time derivative.
The pressure recurrence
applies the incompressibility correction at precisely the same temporal
depth.  Omitting either formula breaks the identity of the differentiated
history.

Spatial differentiation opens an independent direction.  For a multi-index
\(\alpha\in\Nzero^3\), define
\[
 J_{n,\alpha}:=\partial_t^n\partial_x^\alpha u,
 \qquad
 \Pi_{n,\alpha}:=\partial_t^n\partial_x^\alpha p.
\]
Since \(\partial_t\) commutes with fixed spatial derivatives, a second use of
Leibniz' rule yields
\begin{equation}
\begin{aligned}
 J_{n+1,\alpha}
 &+\sum_{a=0}^{n}\sum_{\beta\leq\alpha}
 \binom na\binom\alpha\beta
 (J_{a,\beta}\cdot\nabla)J_{n-a,\alpha-\beta}
 +\nabla\Pi_{n,\alpha}
 \\
 &=\nu\Delta J_{n,\alpha},
 \qquad \nabla\cdot J_{n,\alpha}=0.
\end{aligned}
\label{eq:mixed-recurrence}
\end{equation}
The temporal index counts response at a fixed spatial point; the multi-index
counts spatial resolution of that response.  The Laplacian moves two spatial
orders, the transport term distributes derivatives between two velocity
factors, and pressure couples the result across space.  This is why the two
indices must remain independent.

At \(t=0\), the same recurrence generates the entire finite initial jet.
Write \(V_n^0=V_n(0)\) and \(Q_n^0=Q_n(0)\).  Beginning with \(V_0^0=u_0\),
\begin{align*}
 Q_n^0
 &=R_iR_j\sum_{a=0}^{n}\binom naV_{a,i}^0V_{n-a,j}^0,\\
 V_{n+1}^0
 &=\nu\Delta V_n^0
 -\sum_{a=0}^{n}\binom na(V_a^0\cdot\nabla)V_{n-a}^0
 -\nabla Q_n^0.
\end{align*}
Sobolev multiplication and boundedness of the Riesz transforms preserve every
finite Sobolev order.  Induction gives
\begin{equation}
 V_n^0,Q_n^0\in H^m(\R^3)
 \qquad\text{for every finite }n,m.
 \label{eq:initial-jet}
\end{equation}
The initial condition fixes every finite jet value.  What is not yet known is
whether the corresponding response rows remain integrable with constants
uniform as the terminal cutoff approaches \(T_*\).

\section{Whole-terminal adjacent rectangles}
\label{sec:rectangles}

\subsection{Why the basic unit has two adjacent rows}

Suppose a Hilbert-valued function \(v\) satisfies
\[
 v\in L^2(0,T_*;H^{m+1}),
 \qquad
 v_t\in L^2(0,T_*;H^{m-1}).
\]
The first row measures one derivative more than the desired trace; the second
measures one derivative less.  Their pairing lands exactly at the middle
order.  For smooth \(v\),
\begin{equation}
 \norm{v(t)}_{H^m}^2-\norm{v(r)}_{H^m}^2
 =2\operatorname{Re}\int_r^t
 \pair{v_t(s)}{v(s)}_{H^{m-1},H^{m+1}}\dd s.
 \label{eq:adjacent-identity}
\end{equation}
Cauchy--Schwarz gives
\begin{equation}
 \left|\norm{v(t)}_{H^m}^2-\norm{v(r)}_{H^m}^2\right|
 \leq
 2\norm{v_t}_{L^2(r,t;H^{m-1})}
   \norm{v}_{L^2(r,t;H^{m+1})}.
 \label{eq:adjacent-estimate}
\end{equation}
Density extends the calculation to the stated spaces.  In particular,
\(v\) has a unique representative in \(C([0,T_*];H^m)\); this is the
standard Lions--Magenes trace mechanism \cite{LionsMagenes}.

The estimate explains the geometry needed at the terminal time.  One isolated
Sobolev row does not carry its own time trace.  An adjacent pair does.  A
finite rectangle stores these pairs for finitely many temporal rows and one
chosen spatial depth.

\begin{definition}[Finite mixed rectangle]
\label{def:rectangle}
For \(N,M\in\Nzero\) and \(0<T<T_*\), set
\begin{equation}
\begin{aligned}
 \mathfrak R_{N,M}(u;T)
 :={}&\sum_{n=0}^{N}
 \norm{V_n}_{L^2(0,T;H^{M+1})}^2
 \\
 &+\sum_{n=0}^{N}
 \norm{V_{n+1}}_{L^2(0,T;H^{M-1})}^2.
\end{aligned}
\label{eq:rectangle-norm}
\end{equation}
For \(M=0\), the lower space is the Bessel-potential space \(H^{-1}\).
The object \(\cR_{N,M}[u_0;T]\) is the displayed velocity block together with
the identities \eqref{eq:pressure-recurrence}--\eqref{eq:mixed-recurrence} in
that finite coordinate range.  Its pressure rows are not assigned an
independent norm: after a target order \(m\) is fixed, any rectangle with
spatial depth at least \(m+2\) gives
\[
 Q_n\in L^1(0,T;H^m),
 \qquad \nabla Q_n\in L^1(0,T;H^{m-1})
 \quad(0\leq n\leq N)
\]
by Sobolev multiplication and Riesz boundedness.  We write
\(\cR_{N,M}[u_0]\) for the same finite tuple on \(I_*=(0,T_*)\) when its
whole-terminal norm is finite.
\end{definition}

The notation does not introduce a new solution.  Every entry of
\(\cR_{N,M}[u_0;T]\) is the corresponding derivative of the maximal history
already fixed in \cref{sec:jet}.

\subsection{The datum-generated terminal estimate}

The closure argument begins at the following datum theorem, the first theorem
of the one-fluid compatible-intersection program \cite{BirnieDatum}.  It is
the only nonstandard input used below.

\Needspace{11\baselineskip}
\medskip
\begin{theorem}[Datum input: whole-terminal compatible rectangles]
\label{thm:rectangles}
Let \((u,p)\) be the maximal history generated by
\(u_0\in\Schwartz_\sigma(\R^3)\), and suppose \(T_*<\infty\).  For every
finite \(N,M\in\Nzero\), there is a constant
\(C_{N,M}=C_{N,M}(u_0,\nu,T_*)<\infty\) such that
\begin{equation}
 \sup_{0<T<T_*}\mathfrak R_{N,M}(u;T)
 \leq C_{N,M}.
 \label{eq:whole-terminal-bound}
\end{equation}
If \(N'\leq N\) and \(M'\leq M\), restriction to the lower temporal rows and
the continuous Sobolev embeddings give
\begin{equation}
 \rho_{N',M'}^{N,M}\cR_{N,M}[u_0]
 =\cR_{N',M'}[u_0].
 \label{eq:rectangle-compatibility}
\end{equation}
Every shared entry is the same derivative of the same pressure-complete
history.
\end{theorem}

The decisive content of \cref{thm:rectangles} is its rowwise datum estimate:
for every fixed finite \((n,m)\),
\begin{equation}
 V_n\in L^2(0,T_*;H^{m+1}),
 \qquad
 V_{n+1}\in L^2(0,T_*;H^{m-1}),
 \label{eq:rowwise-datum-theorem}
\end{equation}
with a bound depending on \(u_0,\nu,T_*,n,m\) and uniform in the cutoff
\(T<T_*\).  This is the datum-to-terminal payment.  It is not a consequence
of strict-subinterval smoothness or of the finite summation below.

\begin{proposition}[Finite-coordinate realization]
\label{prop:finite-realization}
The rowwise datum estimate \eqref{eq:rowwise-datum-theorem} is equivalent to
the family of finite bounds \eqref{eq:whole-terminal-bound}, and the resulting
finite blocks obey \eqref{eq:rectangle-compatibility}.
\end{proposition}

\begin{proof}
Monotone convergence in the terminal cutoff turns each terminal-uniform
rowwise bound into the corresponding \(L^2(0,T_*;\cdot)\) membership.
Fix \((N,M)\).  Adding the finitely many squared adjacent-row bounds for
\(0\leq n\leq N\) gives \eqref{eq:whole-terminal-bound}.  Conversely, the
collection of \eqref{eq:whole-terminal-bound} over all finite \((N,M)\)
recovers every instance of \eqref{eq:rowwise-datum-theorem}.  The recurrences
\eqref{eq:pressure-recurrence}--\eqref{eq:mixed-recurrence} identify all
entries with actual derivatives of the one maximal history.  Restricting a
larger finite block leaves each shared derivative unchanged, which proves
\eqref{eq:rectangle-compatibility}.
\end{proof}

Three quantifiers in \cref{thm:rectangles} carry the argument.  The coordinate
\((N,M)\) is fixed before its constant is chosen.  The same constant then
controls every terminal cutoff \(T<T_*\).  No constant independent of
\((N,M)\) occurs.  Thus the theorem provides every finite view of one history
without asking an infinite sum of Sobolev norms to converge.

The distinction from strict-subinterval regularity can be written in one
line:
\[
 \underbrace{\forall T<T_*\ \mathfrak R_{N,M}(u;T)<\infty}_{
 \text{classicality before the endpoint}}
 \quad\not\Rightarrow\quad
 \underbrace{\sup_{0<T<T_*}\mathfrak R_{N,M}(u;T)<\infty}_{
 \text{whole-terminal control}}.
\]
The theorem supplies the expression on the right.  This is what allows the
adjacent trace estimate to see \(T_*\).

\section{The complete zero row reaches one endpoint}
\label{sec:intersection}

Increasing \(N\) reveals more temporal responses and increasing \(M\) reveals
more spatial resolution.  Every overlap is literal equality of derivatives
of \(u\).  In particular, the zeroth temporal entry of the family gives
\begin{equation}
 u\in\Szero(0,T_*):=
 \bigcap_{m=0}^{\infty}L^2(0,T_*;H^m(\R^3)).
 \label{eq:complete-spatial-row}
\end{equation}
For each fixed \(m\), this is one finite membership supplied by a rectangle
of sufficient spatial depth.  The intersection records all such memberships
of the same field; it does not sum their norms.

Write the pressure-completed interaction as
\begin{equation}
 \mathcal B(u,u)
 :=(u\cdot\nabla)u+\nabla p
 =\Leray\nabla\cdot(u\otimes u).
 \label{eq:B}
\end{equation}
The Sobolev product theorem and boundedness of the Leray projection give, for
each integer \(m\geq0\),
\begin{equation}
 \norm{\mathcal B(u,u)}_{H^m}
 \leq C_m\norm{u\otimes u}_{H^{m+1}}
 \leq C_m\norm{u}_{H^{m+2}}^2.
 \label{eq:nonlinear-estimate}
\end{equation}
Standard product and multiplier facts used here are recalled in
\cref{app:analytic}; see also \cite{BCD,Stein}.

Assume for contradiction that \(T_*<\infty\).  From
\eqref{eq:complete-spatial-row} and \eqref{eq:nonlinear-estimate},
\begin{equation}
 \int_0^{T_*}\norm{\mathcal B(u,u)(t)}_{H^m}\dd t
 \leq C_m\norm{u}_{L^2(0,T_*;H^{m+2})}^2<\infty.
 \label{eq:B-L1}
\end{equation}
The viscous term obeys
\begin{equation}
 \int_0^{T_*}\norm{\Delta u(t)}_{H^m}\dd t
 \leq \sqrt{T_*}\norm{u}_{L^2(0,T_*;H^{m+2})}<\infty.
 \label{eq:Delta-L1}
\end{equation}
The original equation, now written as
\[
 u_t=\nu\Delta u-\mathcal B(u,u),
\]
gives
\begin{equation}
 u_t\in L^1(0,T_*;H^m)
 \qquad\text{for every finite }m.
 \label{eq:ut-L1}
\end{equation}
Since \(u\in L^2(0,T_*;H^m)\subset L^1(0,T_*;H^m)\), we have
\begin{equation}
 u\in W^{1,1}(0,T_*;H^m)
 \hookrightarrow C([0,T_*];H^m).
 \label{eq:W11-endpoint}
\end{equation}
More explicitly,
\begin{equation}
\begin{aligned}
 \sup_{0\leq t\leq T_*}\norm{u(t)}_{H^m}
 \leq{}&\norm{u_0}_{H^m}
 +\nu\sqrt{T_*}\norm{u}_{L^2(0,T_*;H^{m+2})}
 \\
 &+C_m\norm{u}_{L^2(0,T_*;H^{m+2})}^2.
\end{aligned}
\label{eq:endpoint-bound}
\end{equation}
Every quantity on the right is finite at its fixed spatial order.  If
\(m_2>m_1\), the two continuous representatives agree almost everywhere in
\(H^{m_1}\), hence everywhere.  They define one distribution
\begin{equation}
 u_*:=u(T_*)\in\bigcap_{m=0}^{\infty}H^m(\R^3)=H^\infty(\R^3).
 \label{eq:endpoint}
\end{equation}
For any sequence \(t_j\uparrow T_*\), the same velocity \(u(t_j)\) converges
to \(u_*\) in every fixed \(H^m\).

The other temporal rows reach this same time by their adjacent pairs.

\medskip
\begin{lemma}[Compatible traces of every response]
\label{lem:all-traces}
For every finite \(n,m\), the response \(V_n\) has a unique representative in
\begin{equation}
 V_n\in C([0,T_*];H^m(\R^3)).
 \label{eq:all-traces}
\end{equation}
Representatives at different Sobolev orders agree after the natural embedding.
For \(n=0\), this representative is the one constructed in
\eqref{eq:W11-endpoint} and its terminal value is \(u_*\).
\end{lemma}

\begin{proof}
Choose a rectangle containing
\[
 V_n\in L^2(0,T_*;H^{m+1}),
 \qquad
 \partial_tV_n=V_{n+1}\in L^2(0,T_*;H^{m-1}).
\]
The adjacent trace estimate \eqref{eq:adjacent-estimate} gives
\eqref{eq:all-traces}.  Two representatives at different orders equal the
same distribution almost everywhere.  Their difference is continuous in the
lower space and vanishes almost everywhere, so it vanishes at every time.
The \(n=0\) representative also equals the \(W^{1,1}\) representative almost
everywhere and hence everywhere.
\end{proof}

\section{The endpoint carries the actual pressure-complete jet}
\label{sec:endpoint-jet}

An \(H^\infty\) velocity endpoint is enough to invoke standard local theory.
The stronger statement needed to preserve the differentiated history is also
available: every left temporal derivative and its normalized pressure reaches
the endpoint, and the limiting values satisfy the same recurrence.

For \(n\in\Nzero\), set
\[
 v_n:=V_n(T_*).
\]
These values exist in every \(H^m\) by \cref{lem:all-traces}.  Define the
normalized endpoint pressure from the velocity traces by
\begin{equation}
 q_n
 =R_iR_j\sum_{a=0}^{n}\binom na v_{a,i}v_{n-a,j}.
 \label{eq:endpoint-pressure}
\end{equation}
The endpoint momentum identity to be verified is
\begin{equation}
 v_{n+1}
 =\nu\Delta v_n
 -\sum_{a=0}^{n}\binom na(v_a\cdot\nabla)v_{n-a}
 -\nabla q_n.
 \label{eq:endpoint-velocity}
\end{equation}

\begin{proposition}[Endpoint identification]
\label{prop:endpoint-jet}
For every finite \(n,m\),
\[
 v_n,q_n\in H^m(\R^3),
\]
and
\begin{equation}
 v_n=\lim_{t\uparrow T_*}\partial_t^n u(t),
 \qquad
 q_n=\lim_{t\uparrow T_*}\partial_t^n p(t)
 \quad\text{in }H^m.
 \label{eq:actual-endpoint-jet}
\end{equation}
Thus \((v_n,q_n)_{n\geq0}\) is the actual left endpoint jet of the maximal
history.
\end{proposition}

\begin{proof}
Every \(v_n\) belongs to \(H^\infty\) by \cref{lem:all-traces}.  Sobolev
multiplication and Riesz boundedness place the right side of
\eqref{eq:endpoint-pressure} in every \(H^m\), so \(q_n\in H^\infty\).
On \((0,T_*)\), the expressions in
\eqref{eq:pressure-recurrence}--\eqref{eq:velocity-recurrence} equal the actual
derivatives \(Q_n\) and \(V_{n+1}\).  Every velocity factor is continuous to
\(T_*\) in arbitrarily high Sobolev spaces by \cref{lem:all-traces}.  The
product estimates and Riesz boundedness allow passage to the endpoint in each
fixed \(H^m\).  The pressure limit is \eqref{eq:endpoint-pressure}; the
momentum limit equals both \(v_{n+1}\), by continuity of \(V_{n+1}\), and the
right side of \eqref{eq:endpoint-velocity}.  This proves the endpoint
identity and both limits in \eqref{eq:actual-endpoint-jet}.
\end{proof}

The pressure normalization matters here.  Adding a different time-dependent
constant to each preterminal pressure would create a formal ambiguity in the
scalar pressure values, even though its gradient is unchanged.  The decaying
Riesz representation \eqref{eq:pressure} fixes one pressure at every time and
passes that same normalization through the endpoint recurrence.

\section{Restart of the same history}
\label{sec:restart}

We now close the contradiction at the single Sobolev level \(H^3\).  The
endpoint \(u_*\in H^\infty_\sigma\) lies in \(H^3_\sigma\), and
\cref{eq:endpoint-bound} gives
\[
 K_3:=\sup_{0\leq t\leq T_*}\norm{u(t)}_{H^3}<\infty.
\]
Standard local well-posedness gives a lifespan
\(\delta=\delta(\nu,K_3)>0\) for every datum in the \(H^3\)-ball of radius
\(K_3\).  Choose \(t_0<T_*\) so close to the endpoint that
\(T_*-t_0<\delta/2\).  The local solution launched from \(u(t_0)\) exists at
least until \(t_0+\delta>T_*\).  On \([t_0,T_*)\) it shares its datum with
the original maximal history, so strong uniqueness identifies the two
solutions there.  It is an extension of that history beyond \(T_*\).

The embedding \(H^3(\R^3)\hookrightarrow W^{1,\infty}(\R^3)\) closes the
transport estimate and strong uniqueness in this local theorem.  Because the
extending solution equals \(u\) on the open overlap, all of their velocity
derivatives and decaying pressure responses agree there.  Continuity to
\(T_*\), together with \cref{prop:endpoint-jet}, identifies the right
endpoint jet of the extension with the original left endpoint jet.  The
continuation is one smooth pressure-complete history, not only an \(H^3\)
extension.

This extension contradicts the definition of \(T_*\) as the maximal classical
time.  The assumption \(T_*<\infty\) is impossible.  We have proved
\(T_*=\infty\), and Sobolev embedding at each finite time gives
\[
 u,p\in C^\infty([0,\infty)\times\R^3).
\]
The energy identity \eqref{eq:energy} remains valid on every finite interval.
Together with the Schwartz datum, divergence constraint, zero forcing, and
decaying pressure normalization fixed at the outset, this supplies the
whole-space finite-energy formulation.  This completes the proof of
\cref{thm:global}.

\section{The critical-height block inside the same family}
\label{sec:critical}

The global contradiction used the complete family directly.  Its familiar
critical quantities can now be read without turning one selected height into
a producer condition.  Put
\[
 \Lambda=(-\Delta)^{1/2},
 \qquad
 E_s(t):=\frac12\norm{\Lambda^s u(t)}_2^2,
 \qquad
 H(t):=E_{1/2}(t).
\]
Pairing \eqref{eq:NS} with \(\Lambda^{2s}u\) gives
\begin{equation}
 E_s'=-2\nu E_{s+1}+\mathcal P_s,
 \label{eq:Sobolev-balance}
\end{equation}
where
\begin{equation}
 \mathcal P_s
 :=-\operatorname{Re}\pair{\Lambda^su}
 {\Lambda^s\bigl((u\cdot\nabla)u+\nabla p\bigr)}.
 \label{eq:production}
\end{equation}
The pressure gradient cancels in this scalar pairing, while its whole-field
action remains in the velocity--pressure recurrence.

\begin{corollary}[Adjacent critical readout]
\label{cor:critical}
The rectangle family gives
\begin{equation}
 E_{3/2}\in L^\infty(0,T_*)\cap L^1(0,T_*),
 \qquad
 H\in W^{1,1}(0,T_*),
 \qquad
 \mathcal P_{1/2}\in L^1(0,T_*).
 \label{eq:critical-block}
\end{equation}
\end{corollary}

\begin{proof}
The rectangle at integer spatial depth \(M=2\) contains the stronger rows
\[
 u\in L^2(0,T_*;H^3),
 \qquad
 u_t\in L^2(0,T_*;H^1).
\]
The trace estimate first gives \(u\in C([0,T_*];H^2)\), hence
\(u\in C([0,T_*];H^{3/2})\) and
\(E_{3/2}\in L^\infty(0,T_*)\).  Its \(L^1\) membership follows from
\(u\in L^2(0,T_*;H^3)\).  The same rows give
\[
 H'(t)=\operatorname{Re}
 \pair{\Lambda^{1/2}u_t(t)}{\Lambda^{1/2}u(t)}\in L^1(0,T_*),
\]
and \(H\in L^1(0,T_*)\).  Hence \(H\in W^{1,1}(0,T_*)\).  At
\(s=1/2\), \eqref{eq:Sobolev-balance} reads
\[
 \mathcal P_{1/2}=H'+2\nu E_{3/2},
\]
whose right side lies in \(L^1(0,T_*)\).
\end{proof}

\appendix

\section{Analytic facts used by the proof}
\label{app:analytic}

We collect the standard estimates used above, both to fix conventions and to
make their finite-order dependence explicit.

\subsection{Sobolev multiplication}

For \(s>3/2\), \(H^s(\R^3)\) is a Banach algebra.  More generally, for each
integer \(m\geq0\), the Leibniz and Sobolev embedding estimates give
\begin{equation}
 \norm{fg}_{H^{m+1}}
 \leq C_m\norm{f}_{H^{m+2}}\norm{g}_{H^{m+2}}.
 \label{eq:product-appendix}
\end{equation}
Applying \(\nabla\cdot\) loses one spatial derivative.  Since the Leray
projection is a matrix of the identity and compositions of Riesz transforms,
it is bounded on every \(H^m\).  Formula \eqref{eq:product-appendix} gives
\eqref{eq:nonlinear-estimate}.  These facts are standard; see
\cite{BCD,Stein}.

\Needspace{18\baselineskip}
\subsection{The adjacent Hilbert triple}

Let
\[
 X_1=H^{m+1}(\R^3),
 \qquad X_0=H^m(\R^3),
 \qquad X_{-1}=H^{m-1}(\R^3).
\]
The dense continuous embeddings
\(X_1\hookrightarrow X_0\hookrightarrow X_{-1}\) identify \(X_{-1}\) with
the dual of \(X_1\) through the pivot \(X_0\).  If
\(v\in L^2(0,T;X_1)\) and \(v_t\in L^2(0,T;X_{-1})\), smooth temporal
convolution gives approximants \(v_\varepsilon\) for which
\eqref{eq:adjacent-identity} is classical.  The two \(L^2\) convergences pass
the right side to the limit, while \eqref{eq:adjacent-estimate} makes the
\(X_0\)-norms Cauchy uniformly at Lebesgue times.  The limit is the unique
\(X_0\)-continuous representative of \(v\).  This is the form of the
Lions--Magenes lemma used in \cref{lem:all-traces}.

\subsection{Spatial smoothness from all finite orders}

For each integer \(k\geq0\), choose \(m>k+3/2\).  The Sobolev embedding
\[
 H^m(\R^3)\hookrightarrow C_b^k(\R^3)
\]
shows
\[
 \bigcap_{m<\infty}H^m(\R^3)\hookrightarrow C^\infty(\R^3).
\]
The constants depend on \(k\) and \(m\).  No uniform constant over all
derivative orders is involved.

\begin{thebibliography}{99}
\small
\setlength{\itemsep}{3pt}

\bibitem{Fefferman}
C. Fefferman,
\href{https://www.claymath.org/wp-content/uploads/2022/06/navierstokes.pdf}
{\emph{Existence and Smoothness of the Navier--Stokes Equation}},
in J. Carlson, A. Jaffe, and A. Wiles (eds.),
\emph{The Millennium Prize Problems}, American Mathematical Society and
Clay Mathematics Institute, 2006, pp.~57--67.

\bibitem{FujitaKato}
H. Fujita and T. Kato,
\emph{On the Navier--Stokes initial value problem. I},
Arch. Rational Mech. Anal. \textbf{16} (1964), 269--315,
\href{https://doi.org/10.1007/BF00276188}{doi:10.1007/BF00276188}.

\bibitem{Kato}
T. Kato,
\emph{Strong \(L^p\)-solutions of the Navier--Stokes equation in
\(\R^m\), with applications to weak solutions},
Math. Z. \textbf{187} (1984), 471--480,
\href{https://doi.org/10.1007/BF01174182}{doi:10.1007/BF01174182}.

\bibitem{BCD}
H. Bahouri, J.-Y. Chemin, and R. Danchin,
\emph{Fourier Analysis and Nonlinear Partial Differential Equations},
Grundlehren der mathematischen Wissenschaften 343, Springer, 2011,
\href{https://doi.org/10.1007/978-3-642-16830-7}
{doi:10.1007/978-3-642-16830-7}.

\bibitem{LionsMagenes}
J.-L. Lions and E. Magenes,
\emph{Non-Homogeneous Boundary Value Problems and Applications, Vol. I},
Springer, 1972, Ch.~III, Lemma~1.2,
\href{https://doi.org/10.1007/978-3-642-65161-8}
{doi:10.1007/978-3-642-65161-8}.

\bibitem{Stein}
E. M. Stein,
\emph{Singular Integrals and Differentiability Properties of Functions},
Princeton Mathematical Series 30, Princeton University Press, 1970.

\bibitem{BirnieDatum}
T. Birnie,
\emph{One-fluid reversible-intersection proof spine},
unpublished manuscript, 2026, Section~5.

\end{thebibliography}

\end{document}
